\documentclass[11pt]{article}

\usepackage[margin=1in]{geometry}
\usepackage{amsmath,amssymb,amsthm,mathtools}
\usepackage{booktabs}
\usepackage[hidelinks,hypertexnames=false]{hyperref}
\usepackage{microtype}

\newtheorem{theorem}{Theorem}[section]
\newtheorem{corollary}[theorem]{Corollary}
\newtheorem{proposition}[theorem]{Proposition}
\newtheorem{lemma}[theorem]{Lemma}
\newtheorem{problem}[theorem]{Open problem}
\theoremstyle{definition}
\newtheorem{definition}[theorem]{Definition}
\theoremstyle{remark}
\newtheorem{remark}[theorem]{Remark}

\newcommand{\N}{\mathbb N}
\newcommand{\R}{\mathbb R}
\newcommand{\C}{\mathbb C}
\newcommand{\PP}{\mathbb P}
\newcommand{\cB}{\mathcal B}
\newcommand{\cG}{\mathcal G}
\newcommand{\cJ}{\mathcal J}
\newcommand{\cS}{\mathcal S}
\newcommand{\diag}{\operatorname{diag}}
\newcommand{\tr}{\operatorname{tr}}
\newcommand{\rank}{\operatorname{rank}}
\newcommand{\indm}{\operatorname{ind}_{-}}
\newcommand{\nm}{n_{-}}
\newcommand{\Rea}{\operatorname{Re}}
\newcommand{\ip}[2]{\langle #1,#2\rangle}
\newcommand{\norm}[1]{\lVert #1\rVert}

\title{The Prime--Stein Compactification Is a Universal Gauge:\\
A No-Go Theorem for a Proposed Route\\
to the Riemann Hypothesis}
\author{}
\date{Draft of July 22, 2026}

\begin{document}
\maketitle

\begin{abstract}
We audit a proposed reciprocal-prime, rank-two Stein route to the Riemann hypothesis and prove
that its apparent new structure is a universal gauge.  For a prime $p$, set
\[
 y_p=\frac{p+1}{p-1},\qquad
 s_p=\frac12+y_p,\qquad
 \lambda_p=\frac{\xi'(s_p)}{\xi(s_p)}
\]
and define
\[
 \cB_{pr}
 =\frac{(p-1)(r-1)}{(p+1)(r+1)(pr-1)}
   \bigl(\lambda_p+\lambda_r\bigr).
\]
Every $s_p$ lies in $(3/2,7/2]$, so the entries use only the absolutely convergent
Euler--Gamma formula for $\xi'/\xi$ in $\Rea s>1$.  The representation identities
\[
 \nm(\cB)
 =\sum_{\substack{\xi(\rho)=0\\ \Rea\rho>1/2}}m(\rho),
 \qquad
 \mathrm{RH}\Longleftrightarrow \cB\geq0.
\]
and, for $\mathsf Q=\diag(1/p)$,
\[
 \cB-\mathsf Q\cB\mathsf Q=av^*+va^*,
\qquad
 v_p=\frac{p-1}{p(p+1)},\quad a_p=\lambda_pv_p.
\]
are correct, but they do not open a new route.  We prove the universal law
\[
 K_{ij}=\frac{a_i a_j}{1-q_iq_j}(\lambda_i+\lambda_j)
 =c_i c_j\frac{\lambda_i+\lambda_j}{x_i+x_j},
 \quad
 x_i=\frac{1+q_i}{1-q_i},\quad
 c_i=\frac{\sqrt2a_i}{1-q_i},
\]
valid for arbitrary disk nodes, amplitudes, and responses.  Rank-two Stein displacement follows
identically.  The prime matrix is therefore a positive diagonal congruence of the classical
logarithmic-derivative Cauchy kernel.  We further show that every countable Hermitian kernel admits
a strictly positive diagonal damping that makes it trace class while preserving the inertia of
every finite principal matrix.  Cyclicity, trace-classness, and the Fredholm determinant therefore
add no arithmetic sign information.  An honest attempt to prove positivity through the Stein
equation, Euler-factor decomposition, and low displacement rank returns exactly the pre-existing
positive-real kernel criterion equivalent to RH.  The proposed route is retired.
\end{abstract}

\paragraph{Logical status.}
Every result labelled theorem, proposition, lemma, or corollary is unconditional.
Problem~\ref{prob:prime-stein} is equivalent to RH and is not proved.  The analytic
negative-index input is classical generalized Nevanlinna theory.  This paper withdraws the claim
that the Prime--Stein representation is a new RH route.  Its positive survivor is a universal
screening theorem for diagonal trace-ideal gauges.

\section{Main results}
\label{sec:intro}

Many RH criteria produce a family of forms, moments, or kernels.  A compact operator with
explicit entries and a low-dimensional displacement boundary can look more rigid than the
underlying kernel.  The purpose of this audit is to test whether that extra structure survives
removal of invertible diagonal gauges.  It does not in the Prime--Stein construction.

Let $\PP$ be the set of rational primes.  Put
\begin{equation}
 q_p=\frac1p,\qquad
 y_p=\frac{1+q_p}{1-q_p}=\frac{p+1}{p-1},\qquad
 s_p=\frac12+y_p=\frac{3p+1}{2(p-1)}.                 \tag{1.1}
\end{equation}
The points $y_p$ are Cayley images of reciprocal primes.  They are distinct and converge to the
interior point $1$ of $(1/2,\infty)$.  Define
\begin{equation}
 \lambda_p=\frac{\xi'(s_p)}{\xi(s_p)},                  \tag{1.2}
\end{equation}
where
\[
 \xi(s)=\frac12s(s-1)\pi^{-s/2}\Gamma(s/2)\zeta(s).
\]
Since $3/2<s_p\leq7/2$, no critical-strip evaluation is needed.

Our operator is the real symmetric prime matrix
\begin{equation}
 \boxed{\quad
 \cB_{pr}=g_{pr}(\lambda_p+\lambda_r),\qquad
 g_{pr}=\frac{(p-1)(r-1)}
 {(p+1)(r+1)(pr-1)}.
 \quad}                                                   \tag{1.3}
\end{equation}

The zeta operator is one value of a universal transform.  Let $U$ be a connected complex
neighborhood of $[1,3]$, and let $\mathcal M_{\R}(U)$ be the multiplicative group of holomorphic
zero-free functions on $U$ that are real and nonzero on $U\cap\R$.  Quotient by nonzero real
constants.  For $F\in\mathcal M_{\R}(U)$ put
\begin{equation}
 \ell_p(F)=\frac{F'(y_p)}{F(y_p)},\qquad
 \mathfrak S_U(F)_{pr}=g_{pr}\bigl(\ell_p(F)+\ell_r(F)\bigr).
                                                               \tag{PSL-0}
\end{equation}

\begin{theorem}[Faithful Prime--Stein logarithm]
\label{thm:logarithm}
The assignment
\[
 \mathfrak S_U:\mathcal M_{\R}(U)/\R^\times
 \longrightarrow \mathcal S_1^{\rm sa}(\ell^2(\PP))
\]
is an injective group homomorphism from multiplication to addition:
\begin{equation}
 \mathfrak S_U(FH)=\mathfrak S_U(F)+\mathfrak S_U(H).       \tag{PSL-1}
\end{equation}
Every image operator satisfies
\begin{equation}
 \mathfrak S_U(F)-\mathsf Q\mathfrak S_U(F)\mathsf Q
 =a_Fv^*+va_F^*,\qquad
 (a_F)_p=\ell_p(F)v_p,                                     \tag{PSL-2}
\end{equation}
with the same cyclic pair $(\mathsf Q,v)$ and displacement rank at most two.

If $F$ extends to a real even Cartwright entire function whose logarithmic-derivative kernel has
minimal generalized Nevanlinna factorization, then
\begin{equation}
 \nm(\mathfrak S_U(F))
 =N_+(F):=\sum_{\substack{F(\alpha)=0\\\Rea\alpha>0}}m(\alpha).
                                                               \tag{PSL-3}
\end{equation}
In particular, for any two such functions $F,H$,
\begin{equation}
 \nm\bigl(\mathfrak S_U(F)+\mathfrak S_U(H)\bigr)
 =\nm(\mathfrak S_U(F))+\nm(\mathfrak S_U(H)).             \tag{PSL-4}
\end{equation}
Thus negative inertia is exactly additive on the Prime--Stein image, despite its failure to be
additive for general self-adjoint trace-class operators.
\end{theorem}

\begin{theorem}[Prime-space inertia law]
\label{thm:inertia}
The matrix in \emph{(1.3)} defines a self-adjoint trace-class operator $\cB$ on
$\ell^2(\PP)$.  Its negative spectral multiplicity is
\begin{equation}
 \nm(\cB)
 =\sum_{\substack{\xi(\rho)=0\\ \Rea\rho>1/2}}m(\rho),      \tag{1.4}
\end{equation}
where both sides may take the value $+\infty$.
\end{theorem}

\begin{corollary}[One-operator RH criterion]
\label{cor:rh}
\begin{equation}
 \mathrm{RH}\quad\Longleftrightarrow\quad \cB\geq0.         \tag{1.5}
\end{equation}
\end{corollary}

\begin{theorem}[Rank-two Prime--Stein law]
\label{thm:stein}
Define
\begin{equation}
 \mathsf Q=\diag(q_p),\qquad
 \mathsf L=\diag(\lambda_p),\qquad
 v_p=\frac{p-1}{p(p+1)},\qquad
 a=\mathsf L v.                                            \tag{1.6}
\end{equation}
Let $\cG=(g_{pr})$.  Then $\cG$ is positive trace class,
\begin{align}
 \cG-\mathsf Q\cG\mathsf Q&=vv^*,                          \tag{1.7}\\
 \cB-\mathsf Q\cB\mathsf Q&=av^*+va^*.                     \tag{1.8}
\end{align}
In particular, the Stein displacement rank of $\cB$ is at most two.  Since
$\norm{\mathsf Q}=1/2$, $\cB$ is the unique trace-class solution of \emph{(1.8)} and
\begin{equation}
 \cB=\sum_{n=0}^{\infty}
 \mathsf Q^n(av^*+va^*)\mathsf Q^n                       \tag{1.9}
\end{equation}
with convergence in trace norm.
\end{theorem}

\begin{theorem}[Minimality]
\label{thm:minimal}
The vector $v$ is cyclic for $\mathsf Q$:
\begin{equation}
 \overline{\operatorname{span}}\{\mathsf Q^nv:n\geq0\}
 =\ell^2(\PP).                                             \tag{1.10}
\end{equation}
Consequently the rank-one realization
\[
 \cG=\sum_{n\geq0}\mathsf Q^nvv^*\mathsf Q^n
\]
is controllable and minimal.  No proper reducing prime subspace containing $v$ carries the
operator $\cG$ or the response law \emph{(1.8)}.
\end{theorem}

\begin{theorem}[Trace-class horizontal spectral flow]
\label{thm:flow}
For $0<\omega<1/2$, set
\begin{equation}
 f_{\omega,p}
 =\frac{\xi(1/2+y_p-\omega)}
        {\xi(1/2+y_p+\omega)},\qquad
 \cB_{\omega,pr}
 =g_{pr}\frac{1-f_{\omega,p}f_{\omega,r}}{2\omega}.         \tag{1.11}
\end{equation}
Then $\cB_\omega$ is self-adjoint trace class.  Outside the countable set of shifts at which the
two shifted Xi functions have a common nonreal zero,
\begin{equation}
 \nm(\cB_\omega)
 =\sum_{\substack{\xi(\rho)=0\\
          \Rea\rho>1/2+\omega}}m(\rho).                    \tag{1.12}
\end{equation}
At an exceptional shift, the right side is read after common-divisor cancellation.  Moreover,
\begin{equation}
 \cB_\omega\longrightarrow\cB
 \quad\text{in trace norm as }\omega\downarrow0.            \tag{1.13}
\end{equation}
\end{theorem}

\begin{theorem}[Fredholm tomography]
\label{thm:fredholm}
The Fredholm determinant
\begin{equation}
 \mathfrak D(t)=\det(I+t\cB)                               \tag{1.14}
\end{equation}
is entire in $t$.  Its positive real zeros, counted with multiplicity, are in bijection with the
negative eigenvalues of $\cB$.  Therefore
\begin{equation}
 N_{(0,\infty)}(\mathfrak D)
 =\sum_{\substack{\xi(\rho)=0\\\Rea\rho>1/2}}m(\rho),       \tag{1.15}
\end{equation}
including infinity, and
\begin{equation}
 \mathrm{RH}
 \quad\Longleftrightarrow\quad
 \mathfrak D(t)>0\ \text{for every }t>0.                   \tag{1.16}
\end{equation}
If $F_N$ is any increasing exhaustion of the primes and
$\mathfrak D_N(t)=\det(I+tP_{F_N}\cB P_{F_N})$, then
$\mathfrak D_N\to\mathfrak D$ locally uniformly on $\C$.
\end{theorem}

Together these results give the proposed law:
\[
 \boxed{
 \begin{gathered}
 \text{off-critical divisor}\\
 =\text{negative spectrum of a minimal prime trace-class state}\\
 =\text{positive-zero divisor of a source Fredholm determinant},\\
 \rank(\cB-\mathsf Q\cB\mathsf Q)\leq2.
 \end{gathered}}                                           \tag{1.17}
\]

\section{Classical analytic detector}
\label{sec:detector}

Set
\begin{equation}
 X(z)=\xi\!\left(\frac12+z\right),\qquad
 L(z)=\frac{X'(z)}{X(z)}.                                  \tag{2.1}
\end{equation}
The function $X$ is real entire, even, and of Cartwright class.  For real $x,y>1/2$, define
\begin{equation}
 J(x,y)=\frac{L(x)+L(y)}{x+y}.                             \tag{2.2}
\end{equation}

\begin{theorem}[Logarithmic-derivative negative index]
\label{thm:Jindex}
The Hermitian kernel $J$ has
\begin{equation}
 \indm(J)
 =\sum_{\substack{\xi(\rho)=0\\ \Rea\rho>1/2}}m(\rho),      \tag{2.3}
\end{equation}
including the value $+\infty$.
\end{theorem}

\paragraph{Proof.}
Evenness removes the linear exponential from the Hadamard factorization of $X$.  Pairing a zero
parameter $\alpha$ with $-\alpha$ gives
\begin{equation}
 \ell_\alpha(z)=\frac1{z-\alpha}+\frac1{z+\alpha}
 =\frac{2z}{z^2-\alpha^2}.                                 \tag{2.4}
\end{equation}
Its contribution to $J$ is
\begin{equation}
 \frac{\ell_\alpha(x)+\ell_\alpha(y)}{x+y}
 =\frac{2(xy-\alpha^2)}
 {(x^2-\alpha^2)(y^2-\alpha^2)}.                           \tag{2.5}
\end{equation}
If $\alpha=i\gamma$ is imaginary, (2.5) is
\begin{equation}
 2\frac{x}{x^2+\gamma^2}\frac{y}{y^2+\gamma^2}
 +2\frac{\gamma}{x^2+\gamma^2}
       \frac{\gamma}{y^2+\gamma^2},                        \tag{2.6}
\end{equation}
a positive kernel of rank at most two.  For real $\alpha>0$, it is
\begin{equation}
 2\frac{x}{x^2-\alpha^2}\frac{y}{y^2-\alpha^2}
 -2\frac{\alpha}{x^2-\alpha^2}
       \frac{\alpha}{y^2-\alpha^2},                        \tag{2.7}
\end{equation}
with one negative square.  A nonreal off-axis quartet gives the Hermitian polarization of
(2.7), with one negative coordinate for each member in $\Rea z>0$.

The rational resolvent features belonging to distinct zeros are linearly independent on every
nonempty real interval.  The minimal Krein--Langer factorization adds these indices without
cancellation.  Equivalently, this is the generalized Nevanlinna theorem for the logarithmic
derivative of a real entire function, restricted to a uniqueness interval; see
Alpay--Rovnyak~\cite{AlpayRovnyak1999}.  If the right-half divisor is infinite, every finite
subdivisor gives a negative subspace of its own cardinality. \qed

\begin{remark}
Theorem~\ref{thm:Jindex} is not claimed as new.  Suzuki's shifted Xi quotient gives the
corresponding generalized Schur family~\cite{Suzuki2012}; the infinitesimal zero-shift limit is
(2.2).  Generalized Schur and Nevanlinna factorization already identify half-plane poles or zeros
with negative squares.
\end{remark}

\section{Reciprocal-prime carrier}
\label{sec:carrier}

Put
\begin{equation}
 c_p=\frac{\sqrt2}{p+1}.                                   \tag{3.1}
\end{equation}
Direct algebra gives
\begin{equation}
 \frac{c_pc_r}{y_p+y_r}
 =\frac{(p-1)(r-1)}
 {(p+1)(r+1)(pr-1)}
 =g_{pr}.                                                   \tag{3.2}
\end{equation}

\begin{proposition}[Positive trace-class carrier]
\label{prop:G}
The matrix $\cG=(g_{pr})$ defines a positive trace-class operator on $\ell^2(\PP)$, and
\begin{equation}
 \tr\cG=\sum_{p\in\PP}\frac{p-1}{(p+1)^3}<\infty.           \tag{3.3}
\end{equation}
\end{proposition}

\paragraph{Proof.}
In $L^2(0,\infty)$ define $u_p(t)=c_pe^{-y_pt}$.  Then
$\ip{u_p}{u_r}=g_{pr}$ and
\[
 \sum_p\norm{u_p}_2^2
 =\sum_p g_{pp}
 =\sum_p\frac{p-1}{(p+1)^3}
 \leq\sum_{n=2}^{\infty}\frac1{n^2}<\infty.
\]
The synthesis map with columns $u_p$ is Hilbert--Schmidt.  Its Gram operator is positive trace
class and has trace (3.3); see Simon~\cite{Simon2005}. \qed

\begin{proposition}[Diagonal congruence]
\label{prop:congruence}
For every finite set $F\subset\PP$,
\begin{equation}
 (\cB_{pr})_{p,r\in F}
 =\diag(c_p)_{p\in F}\,
   \bigl(J(y_p,y_r)\bigr)_{p,r\in F}\,
   \diag(c_p)_{p\in F}.                                    \tag{3.4}
\end{equation}
\end{proposition}

\paragraph{Proof.}
By (1.1), $s_p=1/2+y_p$, so $\lambda_p=L(y_p)$.  Substitute this and (3.2) into (1.3). \qed

The weights $c_p$ are nonzero.  Hence (3.4) preserves negative inertia of every finite
compression.  Their decay makes the infinite realization trace class without losing finite
inertia.

\section{Proof of the inertia law}
\label{sec:proof}

\begin{lemma}[Prime uniqueness]
\label{lem:unique}
Restricting $J$ from $(1/2,\infty)$ to $\{y_p:p\in\PP\}$ preserves its number of negative
squares.
\end{lemma}

\paragraph{Proof.}
The nodes are distinct and $y_p\to1$, an interior point of the zero-free interval.  A meromorphic
kernel-space function vanishing at every $y_p$ vanishes identically.  The corresponding kernel
sections have dense span in the reproducing-kernel Pontryagin space.  A finite-dimensional
negative subspace survives sufficiently accurate approximation of a basis by this span.  Apply
this to negative subspaces of every finite dimension. \qed

\begin{lemma}[Coordinate exhaustion]
\label{lem:compact}
If $A$ is compact self-adjoint on $\ell^2(\PP)$, then
\begin{equation}
 \nm(A)=\sup_{F\Subset\PP}\nm(P_FAP_F).                    \tag{4.1}
\end{equation}
\end{lemma}

\paragraph{Proof.}
Every negative subspace of a compression is negative for $A$.  Conversely, let $M$ be a
finite-dimensional negative subspace.  Compactness of its unit sphere gives $\delta>0$ with
\[
 \ip{Ax}{x}\leq-\delta\norm{x}^2\qquad(x\in M).
\]
Finite coordinate projections converge uniformly to the identity on this sphere.  For a large
finite $F$, $P_FM$ has the same dimension and remains negative. \qed

\paragraph{Proof of Theorem~\ref{thm:inertia}.}
The sequence $(\lambda_p)$ is bounded because $L$ is continuous on the compact zero-free interval
$[1,3]$ in the centered variable.  Thus
\begin{equation}
 \cB=\mathsf L\cG+\cG\mathsf L                             \tag{4.2}
\end{equation}
is self-adjoint trace class by Proposition~\ref{prop:G}.  Lemma~\ref{lem:compact} and
Proposition~\ref{prop:congruence} give
\[
 \nm(\cB)
 =\sup_{F\Subset\PP}\nm\bigl(J(y_p,y_r)\bigr)_{p,r\in F}.
\]
Lemma~\ref{lem:unique} identifies the supremum with $\indm(J)$, and
Theorem~\ref{thm:Jindex} gives (1.4). \qed

\paragraph{Proof of Corollary~\ref{cor:rh}.}
The functional equation reflects every off-critical zero on the left to one on the right.
Therefore RH is equivalent to the right side of (1.4) being zero, hence to $\cB$ having no
negative spectrum. \qed

\section{Stein structure and minimality}
\label{sec:stein}

With $q_p=1/p$ and $v$ as in (1.6),
\begin{equation}
 (1-q_pq_r)g_{pr}=v_pv_r.                                 \tag{5.1}
\end{equation}
This proves (1.7) entrywise.  Since $\mathsf L$ commutes with $\mathsf Q$, (4.2) gives
\[
 \begin{aligned}
 \cB-\mathsf Q\cB\mathsf Q
 &=\mathsf L(\cG-\mathsf Q\cG\mathsf Q)
   +(\cG-\mathsf Q\cG\mathsf Q)\mathsf L\\
 &=(\mathsf Lv)v^*+v(\mathsf Lv)^*
 =av^*+va^*.
 \end{aligned}
\]
On the trace ideal $\cS_1$, $T(Y)=\mathsf QY\mathsf Q$ has norm at most $1/4$.  Thus
$I-T$ is invertible by a Neumann series, proving (1.9) and uniqueness.

\paragraph{Proof of Theorem~\ref{thm:minimal}.}
Suppose $x\in\ell^2(\PP)$ is orthogonal to every $\mathsf Q^nv$.  Since both $x$ and $v$ are in
$\ell^2$, the signed measure
\[
 \mu_x=\sum_{p\in\PP}x_pv_p\,\delta_{q_p}
\]
has finite total variation on $[0,1/2]$.  Orthogonality says every moment of $\mu_x$ vanishes.
Polynomials are uniformly dense in $C[0,1/2]$, so $\mu_x=0$.  Each $v_p$ is nonzero; hence every
$x_p=0$.  This proves cyclicity.  Minimality follows from (1.7) and the standard controllability
characterization. \qed

\begin{corollary}[All-depth prime correlation]
\label{cor:depth}
For every finitely supported real vector $u=(u_p)$,
\begin{equation}
 \ip{\cB u}{u}
 =2\sum_{n=0}^{\infty}A_n(u)V_n(u),                        \tag{5.2}
\end{equation}
where
\begin{equation}
 A_n(u)=\sum_p a_pq_p^nu_p,\qquad
 V_n(u)=\sum_p v_pq_p^nu_p.                                \tag{5.3}
\end{equation}
\end{corollary}

The old source-to-Xi morphism failed to specify the two signals, their shift, their completion,
or their pairing.  Equations (5.2)--(5.3) specify all of them.  The RH demand is the sign of the
complete valuation-depth correlation.

\section{Trace-class aperture flow}
\label{sec:flow}

For $0<\omega<1/2$, define $f_{\omega,p}$ and $\cB_\omega$ by (1.11), and put
$\mathsf F_\omega=\diag(f_{\omega,p})$.  Both Xi arguments in (1.11) exceed $1$, so
$\mathsf F_\omega$ is bounded and source-defined by convergent Euler products.  Equation (3.2)
gives
\begin{equation}
 \cB_\omega
 =\frac{\cG-\mathsf F_\omega\cG\mathsf F_\omega}{2\omega}.
                                                               \tag{6.1}
\end{equation}
Thus $\cB_\omega$ is self-adjoint trace class and
\begin{equation}
 \cB_\omega-\mathsf Q\cB_\omega\mathsf Q
 =\frac{vv^*-(\mathsf F_\omega v)(\mathsf F_\omega v)^*}
        {2\omega},                                         \tag{6.2}
\end{equation}
again of displacement rank at most two.

\paragraph{Proof of the index identity.}
The shifted real-ray kernel is
\begin{equation}
 J_\omega(x,y)
 =\frac{1-f_\omega(x)f_\omega(y)}{2\omega(x+y)},\qquad
 f_\omega(x)=\frac{\xi(1/2+x-\omega)}
                   {\xi(1/2+x+\omega)}.                    \tag{6.3}
\end{equation}
The generalized Schur index theorem gives, away from common-divisor shifts,
\[
 \indm(J_\omega)
 =\sum_{\Rea\rho>1/2+\omega}m(\rho).
\]
The prime set $\{y_p\}$ is still a uniqueness set, and
$\cB_{\omega,pr}=c_pc_rJ_\omega(y_p,y_r)$.  The proof of
Theorem~\ref{thm:inertia} gives (1.12).

\paragraph{Proof of trace-norm convergence.}
Uniformly for $x\in[1,3]$,
\begin{equation}
 f_\omega(x)=1-2\omega L(x)+O(\omega^2).                   \tag{6.4}
\end{equation}
Consequently
\[
 \frac{I-\mathsf F_\omega}{2\omega}\longrightarrow\mathsf L
\quad\text{in operator norm}.
\]
Writing
\[
 \cG-\mathsf F_\omega\cG\mathsf F_\omega
 =(I-\mathsf F_\omega)\cG
 +\mathsf F_\omega\cG(I-\mathsf F_\omega)
\]
and using $\cG\in\cS_1$ proves $\cB_\omega\to
\mathsf L\cG+\cG\mathsf L=\cB$ in trace norm. \qed

\begin{corollary}[Horizontal divisor moments]
\label{cor:moments}
For every $\alpha>0$,
\begin{equation}
 \alpha\int_0^{1/2}\omega^{\alpha-1}\nm(\cB_\omega)\,d\omega
 =\sum_{\substack{\xi(\rho)=0\\\Rea\rho>1/2}}
   m(\rho)(\Rea\rho-1/2)^\alpha,                           \tag{6.5}
\end{equation}
with $+\infty$ allowed.  Countably many exceptional shifts do not change the integral.
\end{corollary}

\paragraph{Proof.}
Outside the exceptional set, Theorem~\ref{thm:flow} identifies the index with the cumulative
right-half zero count.  Apply Tonelli's theorem and
$\alpha\int_0^a\omega^{\alpha-1}d\omega=a^\alpha$ to each zero. \qed

The family contains more than the Boolean RH condition: its negative-index step function recovers
horizontal distances and multiplicities, and its layer-cake moments recover every positive
transverse moment.

\section{Fredholm tomography}
\label{sec:fredholm}

\paragraph{Proof of Theorem~\ref{thm:fredholm}.}
Since $\cB$ is trace class, its Fredholm determinant has the convergent product
\begin{equation}
 \mathfrak D(t)=\prod_j(1+t\mu_j),                          \tag{7.1}
\end{equation}
where $(\mu_j)$ are nonzero eigenvalues of $\cB$, counted with multiplicity.  A positive real zero
$t$ occurs exactly when $\mu_j=-1/t<0$, with the same multiplicity.  Equation (1.15) follows from
Theorem~\ref{thm:inertia}.  Under RH every $\mu_j\geq0$, so every factor is positive for $t>0$;
the converse follows from the positive root attached to a negative eigenvalue.

For an increasing prime exhaustion, $P_{F_N}\cB P_{F_N}\to\cB$ in trace norm.  Continuity of the
Fredholm determinant in trace norm gives locally uniform convergence of $\mathfrak D_N$. \qed

The same construction applies to every aperture:
\begin{equation}
 \mathfrak D_\omega(t)=\det(I+t\cB_\omega).                \tag{7.2}
\end{equation}
Away from cancellation shifts, the positive-zero count of $\mathfrak D_\omega$ equals the number
of zeros in $\Rea\rho>1/2+\omega$.  The bivariate source object
$(\omega,t)\mapsto\mathfrak D_\omega(t)$ is a Fredholm tomography of the transverse divisor.

\section{Euler-source form and open theorem}
\label{sec:open}

For $s>1$,
\begin{equation}
 \frac{\xi'(s)}{\xi(s)}
 =\frac1s+\frac1{s-1}-\frac12\log\pi
   +\frac12\psi(s/2)
   -\sum_{m=2}^{\infty}\frac{\Lambda(m)}{m^s},              \tag{8.1}
\end{equation}
where $\Lambda$ is the von Mangoldt function.  Let
\[
 h(s)=\frac1s+\frac1{s-1}-\frac12\log\pi+\frac12\psi(s/2),
 \quad D_h=\diag(h(s_p)),\quad D_m=\diag(m^{-s_p}).
\]
Then
\begin{equation}
 \cB
 =D_h\cG+\cG D_h
 -\sum_{m=2}^{\infty}\Lambda(m)(D_m\cG+\cG D_m).            \tag{8.2}
\end{equation}
The series converges in trace norm because
\begin{equation}
 \norm{D_m\cG+\cG D_m}_1
 \leq2m^{-3/2}\norm{\cG}_1,\qquad
 \sum_{m\geq2}\Lambda(m)m^{-3/2}<\infty.                   \tag{8.3}
\end{equation}

The Prime--Stein logarithm makes the Euler product itself an operator sum.  In centered variable
$y$, write
\begin{equation}
 X(y)=\xi(1/2+y)=A_\infty(y)\prod_{\ell\in\PP}E_\ell(y),
 \qquad
 E_\ell(y)=\bigl(1-\ell^{-1/2-y}\bigr)^{-1},               \tag{8.4}
\end{equation}
where $A_\infty$ is the polynomial--Gamma completion.

\begin{theorem}[Trace-norm Euler factorization]
\label{thm:euler-factorization}
On a complex neighborhood of the sample interval $[1,3]$ contained in $\Rea y>1/2$,
\begin{equation}
 \boxed{\quad
 \cB=\mathfrak S(A_\infty)
      +\sum_{\ell\in\PP}\mathfrak S(E_\ell),
 \quad}                                                     \tag{8.5}
\end{equation}
and the sum converges absolutely in trace norm.  Every summand has Stein displacement rank at
most two with respect to the same cyclic pair $(\mathsf Q,v)$.
\end{theorem}

\paragraph{Proof.}
Finite Euler products satisfy (8.5) by the homomorphism law (PSL-1).  At a sample point,
\[
 \left|\frac{E_\ell'(y_p)}{E_\ell(y_p)}\right|
 =\frac{(\log\ell)\ell^{-s_p}}{1-\ell^{-s_p}}
 \leq\frac{(\log\ell)\ell^{-3/2}}{1-\ell^{-3/2}}.
\]
The bound is summable over primes.  Since
$\norm{D\cG+\cG D}_1\leq2\norm D\norm{\cG}_1$, the operator series converges absolutely in trace
norm to the logarithmic-derivative response of the full Euler product.  Equation (PSL-2) gives
the common displacement law. \qed

\begin{problem}[Prime--Stein positivity]
\label{prob:prime-stein}
Prove directly from \emph{(8.1)} or \emph{(8.2)}, without using Xi zero locations, that
\begin{equation}
 \cB\geq0.                                                  \tag{PS}
\end{equation}
Equivalently, prove for every finitely supported real prime packet $u$ that
\begin{equation}
 \sum_{p,r\in\PP}
 \frac{(p-1)(r-1)}
 {(p+1)(r+1)(pr-1)}
 (\lambda_p+\lambda_r)u_pu_r\geq0.                         \tag{8.6}
\end{equation}
\end{problem}

Problem~\ref{prob:prime-stein} is exactly RH-strength.  The proposed normal form is:
\begin{equation}
 \begin{gathered}
 \text{fixed strict contraction }\mathsf Q=\diag(1/p),\\
 \text{cyclic boundary vector }v,\quad
 \text{explicit Euler response }a,\quad
 \rank(av^*+va^*)\leq2,\\
 \text{unique trace-class state }
 \cB=(I-\mathsf Q(\,\cdot\,)\mathsf Q)^{-1}(av^*+va^*),\\
 \mathrm{RH}\Longleftrightarrow\cB\geq0.
 \end{gathered}                                             \tag{8.7}
\end{equation}
Section~\ref{sec:nogo} proves that this normal form is a universal Cayley gauge of the classical
positive-real kernel.  It therefore labels the frontier as structured-operator passivity without
moving it.

\paragraph{Finite falsification.}
If RH is false, Theorem~\ref{thm:inertia} and Lemma~\ref{lem:compact} give a finite set of primes
$F$ and a real vector supported on $F$ with negative value in (8.6).  Conversely, any such packet
proves RH false.  Failure has a finite source certificate even though positivity is all-order.

\paragraph{Why positivity is not automatic.}
Both $\cG$ and $\mathsf L$ have positive diagonal data in the Euler region, but the
anticommutator $\mathsf L\cG+\cG\mathsf L$ need not be positive for a general positive $\cG$.
Already at two nodes, unequal positive responses can make a Cauchy--Loewner determinant negative.
The missing sign is a global arithmetic correlation, not scalar monotonicity of $\xi$.

\section{Decisive gauge no-go}
\label{sec:nogo}

\begin{theorem}[Universal Cauchy--Stein gauge]
\label{thm:universal-gauge}
Let $I$ be any index set and choose real data $q_i,a_i,\lambda_i$ such that the denominators below
do not vanish.  Set
\[
 x_i=\frac{1+q_i}{1-q_i},\qquad
 G_{ij}=\frac{a_i a_j}{1-q_iq_j},\qquad
 K_{ij}=G_{ij}(\lambda_i+\lambda_j),\qquad
 c_i=\frac{\sqrt2a_i}{1-q_i}.
\]
Then
\begin{align}
 G_{ij}-q_iq_jG_{ij}&=a_ia_j, \tag{9.1}\\
 K_{ij}-q_iq_jK_{ij}
   &=(\lambda_i a_i)a_j+a_i(\lambda_j a_j), \tag{9.2}\\
 K_{ij}&=c_i c_j\frac{\lambda_i+\lambda_j}{x_i+x_j}. \tag{9.3}
\end{align}
Thus the carrier displacement, response displacement, and Cauchy congruence are universal.
\end{theorem}

\paragraph{Proof.}
The Cayley identity
\[
 \frac1{x_i+x_j}
 =\frac{(1-q_i)(1-q_j)}{2(1-q_iq_j)}
\]
gives (9.3); multiplying the definitions by $1-q_iq_j$ gives (9.1)--(9.2). \qed

For the prime construction take
\[
 q_p=\frac1p,\qquad
 a_p=\frac{q_p(1-q_p)}{1+q_p}.
\]
Then $G_{pr}=g_{pr}$ and $c_p=\sqrt2/(p+1)$.  Hence the diagonal congruence (3.4) is not merely analogous to the
universal theorem: it is its literal specialization.  Primality contributes a convenient
uniqueness sequence and summable weights, but no new sign identity.

\begin{theorem}[Free trace-ideal compactification]
\label{thm:free-gauge}
Every finite-valued Hermitian kernel $H=(h_{mn})$ on a countable set admits a strictly positive
diagonal $D=\diag(d_n)$ such that $DHD$ is trace class and every finite principal matrix of
$DHD$ has the same inertia as the corresponding principal matrix of $H$.
\end{theorem}

\paragraph{Proof.}
Index the set by positive integers and put
\[
 M_n=\max\bigl(1,|h_{ij}|:1\leq i,j\leq n\bigr),
 \qquad d_n=\frac{2^{-n}}{M_n}.
\]
If $i\leq j$, then
\[
 |d_i h_{ij}d_j|\leq2^{-i-j}.
\]
The entries of $DHD$ are therefore absolutely summable, so $DHD$ is trace class.  Each finite
restriction of $D$ is positive and invertible; Sylvester's law of inertia proves the second
claim. \qed

Theorem~\ref{thm:free-gauge} shows why trace-classness and the Fredholm determinant cannot certify
a new arithmetic mechanism.  They can be added to every countable kernel without changing a
single finite sign witness.  Cyclicity does not repair this: distinct nodes and nonzero boundary
amplitudes already give the same weighted-Vandermonde minimality for arbitrary data.

Finally, the prime nodes converge to the interior point $1$.  Analytic uniqueness and locality of
negative squares identify their restricted kernel with the full kernel $J$.  Thus
\[
 \cB\geq0
 \quad\Longleftrightarrow\quad
 J\geq0
 \quad\Longleftrightarrow\quad
 \mathrm{RH}.
\]
The first equivalence is diagonal congruence plus uniqueness; it is not a new antecedent for RH.
The attempted route has returned to its starting criterion.

\section{Universal completed-Euler-product law}
\label{sec:universal}

\paragraph{Proof of Theorem~\ref{thm:logarithm}.}
The sample values $\ell_p(F)$ are bounded because $y_p\to1$ inside $U$.  With
$D_F=\diag(\ell_p(F))$, definition (PSL-0) is
\begin{equation}
 \mathfrak S_U(F)=D_F\cG+\cG D_F,                          \tag{10.1}
\end{equation}
so Proposition~\ref{prop:G} makes it self-adjoint trace class.  The logarithmic derivative turns
products into sums, proving (PSL-1).  Equation (PSL-2) follows from (1.7), because $D_F$ commutes
with $\mathsf Q$.

For faithfulness, suppose $\mathfrak S_U(F)=\mathfrak S_U(H)$.  Since $g_{pp}>0$, equality of
diagonal entries gives
\[
 \ell_p(F)=\ell_p(H)\qquad(p\in\PP).
\]
The two logarithmic derivatives are holomorphic on $U$ and agree on points accumulating at
$1\in U$.  The identity theorem gives $F'/F=H'/H$ throughout $U$, so $(F/H)'=0$ and $F/H$ is a
nonzero real constant.  Hence the quotient map is injective.

If $F$ has the stated real-even Cartwright extension, apply the proof of
Theorem~\ref{thm:inertia} with $X$ replaced by $F$.  The same uniqueness set transfers its minimal
generalized Nevanlinna index to $\mathfrak S_U(F)$, proving (PSL-3).  Finally
\[
 N_+(FH)=N_+(F)+N_+(H)
\]
with multiplicity.  Combine this divisor identity, (PSL-1), and (PSL-3) to obtain the exceptional
inertia additivity (PSL-4). \qed

\begin{corollary}[Completed-Euler-product index law]
\label{cor:universal}
For every self-dual completed Euler product whose centered completion satisfies the Cartwright,
Euler-convergence, and minimal-factorization hypotheses, the Prime--Stein logarithm is a faithful
source operator and its negative index counts failures of the corresponding generalized RH.
Products of such completions become sums of trace-class operators, and their off-critical counts
become additive inertia.
\end{corollary}

Odd functional-equation parity adds a central term $m/z$ to the logarithmic derivative.  Its
kernel is the positive rank-one term $m/(xy)$ and does not alter negative inertia.  The carrier
$(\mathsf Q,v,\cG)$ is universal; each completed Euler product enters only through its boundary
response $a_F$.  Thus factorization of completed $L$-functions is transported faithfully into
addition inside one fixed rank-two-displacement trace ideal.

\section{Novelty and literature boundary}
\label{sec:novelty}

Known ingredients, not claimed as contributions, are:
\begin{itemize}
\item Weil- and Li-type positivity criteria for RH~\cite{Weil1952,Li1997};
\item de Branges and generalized Schur/Nevanlinna negative-square theory
      \cite{deBranges1968,AlpayRovnyak1999};
\item Suzuki's shifted-Xi meromorphic-inner criterion~\cite{Suzuki2012};
\item Cauchy Gram matrices, trace ideals, Stein equations, and Fredholm determinants
      \cite{Aronszajn1950,Simon2005};
\item the logarithmic-derivative kernel in Theorem~\ref{thm:Jindex}.
\end{itemize}

The hostile comparison changes the scientific verdict.  The exact negative-index law is the
classical generalized-Nevanlinna kernel localized to a uniqueness set.  Alpay--Rovnyak's
region-independence of negative squares already supplies the analytic mechanism.  The
reciprocal-prime carrier is the arbitrary-node Schur carrier of
Theorem~\ref{thm:universal-gauge}; trace-classness is freely imposed by
Theorem~\ref{thm:free-gauge}; and the Fredholm zero statement is the spectral theorem after that
gauge.  Searching for the package as a phrase was therefore the wrong novelty test: a change of
coordinates hides the collision with classical parent theory.

The strongest surviving contribution is methodological and exact.  Any proposed RH operator
whose novelty consists only of countable uniqueness sampling, positive diagonal damping,
low Stein displacement rank, cyclicity, or a Fredholm determinant must first be quotiented by
Theorems~\ref{thm:universal-gauge}--\ref{thm:free-gauge}.  Only structure that survives this
quotient can constitute a new route.  This screening principle is useful, but its ingredients are
elementary and classical; we do not claim new scientific priority for it.

\section{Conclusion}

The Prime--Stein formulas produce a minimal trace-class operator from Euler-region values, with
negative spectrum equal to the off-critical divisor and displacement boundary of rank two:
\[
 \boxed{
 \begin{array}{c}
 \displaystyle
 \lambda_p=
 \left.\frac{\xi'}{\xi}\right|_{(3p+1)/(2(p-1))}
 \\[2mm]
 \Downarrow\\[-1mm]
 \cB-\mathsf Q\cB\mathsf Q=av^*+va^*
 \\[1mm]
 \Downarrow\\[-1mm]
 \nm(\cB)=N_{\Rea\rho>1/2}
 \\[1mm]
 \Downarrow\\[-1mm]
 \det(I+t\cB)>0\ \forall t>0
 \Longleftrightarrow\cB\geq0
 \Longleftrightarrow\mathrm{RH}.
 \end{array}}
\]
But every arrow in this display after the first is invariant under the universal diagonal gauge.
The trace-class operator is a compactified copy of the pre-existing Cauchy kernel, and its
positivity is that kernel's RH criterion verbatim.  Solving the Stein equation gives existence,
not passivity; Euler factors have negative logarithmic-derivative diagonal and require the same
global archimedean--prime cancellation; low displacement rank and cyclicity impose no sign.

Accordingly this draft does not present a new route to RH.  It records a failed route and the
exact theorem that kills it.  A successor must introduce an arithmetic constraint that survives
finite positive diagonal congruence and proves part of the three-node-and-higher sign problem
before invoking trace ideals, Stein realizations, or Fredholm determinants.

\begin{thebibliography}{99}

\bibitem{AlpayRovnyak1999}
D.~Alpay and J.~Rovnyak,
\emph{Loewner's theorem for kernels having a finite number of negative squares},
Proceedings of the American Mathematical Society \textbf{127} (1999), 1109--1117,
\href{https://doi.org/10.1090/S0002-9939-99-04618-3}{doi:10.1090/S0002-9939-99-04618-3}.

\bibitem{Aronszajn1950}
N.~Aronszajn,
\emph{Theory of reproducing kernels},
Transactions of the American Mathematical Society \textbf{68} (1950), 337--404.

\bibitem{deBranges1968}
L.~de Branges,
\emph{Hilbert Spaces of Entire Functions},
Prentice--Hall, 1968.

\bibitem{Li1997}
X.-J.~Li,
\emph{The positivity of a sequence of numbers and the Riemann hypothesis},
Journal of Number Theory \textbf{65} (1997), 325--333.

\bibitem{Simon2005}
B.~Simon,
\emph{Trace Ideals and Their Applications}, second edition,
American Mathematical Society, 2005.

\bibitem{Suzuki2012}
M.~Suzuki,
\emph{A canonical system of differential equations arising from the Riemann zeta-function},
arXiv:1204.1827 (2012),
\url{https://arxiv.org/abs/1204.1827}.

\bibitem{Titchmarsh1986}
E.~C.~Titchmarsh, revised by D.~R.~Heath-Brown,
\emph{The Theory of the Riemann Zeta-Function}, second edition,
Oxford University Press, 1986.

\bibitem{Weil1952}
A.~Weil,
\emph{Sur les formules explicites de la th\'eorie des nombres premiers},
Communications du S\'eminaire Math\'ematique de l'Universit\'e de Lund (1952), 252--265.

\end{thebibliography}

\end{document}
